% GENERATED FILE. Edit canonical lessons, then run npm run generate:books.
% canonical-source-hash: fnv1a64:b5a3f8dfc741757a
% canonical-lessons: IT-C351-origine, IT-C351-derivare, IT-C351-comportare, IT-C351-determinare, IT-C351-influire

\chapter{Da dove nasce}
\label{ch:it-da-dove-nasce}
\hlchaptermodality{351}

\begin{chapteropening}
\textbf{By the end of this chapter:} \emph{I can say where a problem comes from and what it involves, determines or influences.}

Complete the last lesson of chapter 351: I can say where a problem comes from and what it involves, determines or influences.
\end{chapteropening}

\section[l'origine]{l'origine — the origin}

\label{lesson:IT-C351-origine}



\begin{quote}
Before the new one: say the Italian for therefore, then the Italian for due.
\end{quote}

\subsection*{You'll want to know: l'origine}
\textbf{l'origine} — \textquotedblleft{}the origin\textquotedblright{}.

\begin{quote}
Qual è l'origine di questa tradizione?
\end{quote}

\begin{quote}
\emph{What is the origin of this tradition?}
\end{quote}

\emph{Avere origine da} — to come from. \emph{Di origine italiana} — of Italian origin.

\begin{cousinweb}
\emph{origine} is a learned borrowing from Latin \emph{orīgō, orīginis} \textquotedblleft{}beginning, source, birth\textquotedblright{} (Lewis \& Short), from \emph{orīrī} \textquotedblleft{}to rise\textquotedblright{}, as the sun rises.
\end{cousinweb}

\begin{grammarlens}[title={What you've built}]
One more word for telling, explaining and describing at B1.
\end{grammarlens}

\subsection*{Your turn}
\begin{itemize}
\raggedright
  \item \emph{Say it:} \emph{l'origine}
  \item \emph{Say it:} \emph{l'origine}, once more
  \item \emph{Say it:} a sentence of your own with \emph{l'origine}
  \item \emph{Recall:} say the Italian for therefore, then the Italian for due, then say \emph{l'origine} again
\end{itemize}

\begin{tcolorbox}[breakable,colback=teal!4,colframe=teal!35!black,title={Before you move on}]
What does l'origine mean? (\textquotedblleft{}The origin\textquotedblright{}.) Say it once more.
\end{tcolorbox}
\section[derivare]{derivare — to come from, to derive (from)}

\label{lesson:IT-C351-derivare}



\begin{quote}
Before the new one: say the Italian for due, then the Italian for the origin.
\end{quote}

\subsection*{You'll want to know: derivare}
\textbf{derivare} — \textquotedblleft{}to come from, to derive (from)\textquotedblright{}.

\begin{quote}
Molti problemi derivano dalla mancanza di soldi.
\end{quote}

\begin{quote}
\emph{Many problems come from a lack of money.}
\end{quote}

\emph{Derivare da} takes \emph{essere}: \emph{il problema è derivato da...}

\begin{cousinweb}
\emph{derivare} is a learned borrowing from Latin \emph{dērīvāre} \textquotedblleft{}to lead off, to draw off water\textquotedblright{} (Lewis \& Short), built from \emph{dē} \textquotedblleft{}from\textquotedblright{} and \emph{rīvus} \textquotedblleft{}a stream\textquotedblright{}: a channel flowing out of a river.
\end{cousinweb}

\begin{grammarlens}[title={What you've built}]
One more word for telling, explaining and describing at B1.
\end{grammarlens}

\subsection*{Your turn}
\begin{itemize}
\raggedright
  \item \emph{Say it:} \emph{derivare}
  \item \emph{Say it:} \emph{derivare}, once more
  \item \emph{Say it:} a sentence of your own with \emph{derivare}
  \item \emph{Recall:} say the Italian for due, then the Italian for the origin, then say \emph{derivare} again
\end{itemize}

\begin{tcolorbox}[breakable,colback=teal!4,colframe=teal!35!black,title={Before you move on}]
What does derivare mean? (\textquotedblleft{}To come from, to derive (from)\textquotedblright{}.) Say it once more.
\end{tcolorbox}
\section[comportare]{comportare — to involve, to entail}

\label{lesson:IT-C351-comportare}



\begin{quote}
Before the new one: say the Italian for the origin, then the Italian for to come from.
\end{quote}

\subsection*{You'll want to know: comportare}
\textbf{comportare} — \textquotedblleft{}to involve, to entail\textquotedblright{}.

\begin{quote}
Questo lavoro comporta molti viaggi.
\end{quote}

\begin{quote}
\emph{This job involves a lot of travel.}
\end{quote}

\emph{Comportarsi} is a different use: to behave. \emph{Comportare delle spese} — to involve costs.

\begin{cousinweb}
\emph{comportare} comes from Latin \emph{comportāre} \textquotedblleft{}to carry together, to bring in\textquotedblright{} (Lewis \& Short), built from \emph{com-} and \emph{portāre} \textquotedblleft{}to carry\textquotedblright{}: what a choice carries along with it.
\end{cousinweb}

\begin{grammarlens}[title={What you've built}]
One more word for telling, explaining and describing at B1.
\end{grammarlens}

\subsection*{Your turn}
\begin{itemize}
\raggedright
  \item \emph{Say it:} \emph{comportare}
  \item \emph{Say it:} \emph{comportare}, once more
  \item \emph{Say it:} a sentence of your own with \emph{comportare}
  \item \emph{Recall:} say the Italian for the origin, then the Italian for to come from, then say \emph{comportare} again
\end{itemize}

\begin{tcolorbox}[breakable,colback=teal!4,colframe=teal!35!black,title={Before you move on}]
What does comportare mean? (\textquotedblleft{}To involve, to entail\textquotedblright{}.) Say it once more.
\end{tcolorbox}
\section[determinare]{determinare — to determine}

\label{lesson:IT-C351-determinare}



\begin{quote}
Before the new one: say the Italian for to come from, then the Italian for to involve.
\end{quote}

\subsection*{You'll want to know: determinare}
\textbf{determinare} — \textquotedblleft{}to determine\textquotedblright{}.

\begin{quote}
Il tempo determina il nostro programma per domenica.
\end{quote}

\begin{quote}
\emph{The weather determines our plan for Sunday.}
\end{quote}

\emph{Determinato} also means determined, resolute: \emph{una ragazza determinata}.

\begin{cousinweb}
\emph{determinare} is a learned borrowing from Latin \emph{dēterminō} \textquotedblleft{}to limit, to set a bound\textquotedblright{} (EtymDB), built on \emph{terminus} \textquotedblleft{}a boundary stone\textquotedblright{}.
\end{cousinweb}

\begin{grammarlens}[title={What you've built}]
One more word for telling, explaining and describing at B1.
\end{grammarlens}

\subsection*{Your turn}
\begin{itemize}
\raggedright
  \item \emph{Say it:} \emph{determinare}
  \item \emph{Say it:} \emph{determinare}, once more
  \item \emph{Say it:} a sentence of your own with \emph{determinare}
  \item \emph{Recall:} say the Italian for to come from, then the Italian for to involve, then say \emph{determinare} again
\end{itemize}

\begin{tcolorbox}[breakable,colback=teal!4,colframe=teal!35!black,title={Before you move on}]
What does determinare mean? (\textquotedblleft{}To determine\textquotedblright{}.) Say it once more.
\end{tcolorbox}
\section[influire]{influire — to influence, to have an effect (on)}

\label{lesson:IT-C351-influire}



\begin{quote}
Before the new one: say the Italian for to involve, then the Italian for to determine.
\end{quote}

\subsection*{You'll want to know: influire}
\textbf{influire} — \textquotedblleft{}to influence, to have an effect (on)\textquotedblright{}.

\begin{quote}
Il tempo influisce sulle nostre scelte.
\end{quote}

\begin{quote}
\emph{The weather has an effect on our choices.}
\end{quote}

\emph{Influire su} qualcosa. Conjugated with \emph{-isc-}: \emph{influisce}.

\begin{cousinweb}
\emph{influire} is a learned borrowing from Latin \emph{influere} \textquotedblleft{}to flow in\textquotedblright{} (Lewis \& Short), built from \emph{in} and \emph{fluere} \textquotedblleft{}to flow\textquotedblright{}; \emph{influenza} comes from the same verb.
\end{cousinweb}

\begin{grammarlens}[title={What you've built}]
One more word for telling, explaining and describing at B1.
\end{grammarlens}

\subsection*{Your turn}
\begin{itemize}
\raggedright
  \item \emph{Say it:} \emph{influire}
  \item \emph{Say it:} \emph{influire}, once more
  \item \emph{Say it:} a sentence of your own with \emph{influire}
  \item \emph{Recall:} say the Italian for to involve, then the Italian for to determine, then say \emph{influire} again
\end{itemize}

\begin{tcolorbox}[breakable,colback=teal!4,colframe=teal!35!black,title={Before you move on}]
What does influire mean? (\textquotedblleft{}To influence, to have an effect (on)\textquotedblright{}.) Say it once more.
\end{tcolorbox}
